\documentclass[12pt]{article}

% --- パッケージ ---
\usepackage[margin=2.5cm]{geometry}
\usepackage{amsmath,amssymb,amsthm}
\usepackage{hyperref}
\usepackage{booktabs}
\usepackage{setspace}
\usepackage{natbib}
\usepackage[T1]{fontenc}


% --- 定理環境 ---
\newtheorem{theorem}{Theorem}[section]
\newtheorem{lemma}[theorem]{Lemma}
\newtheorem{corollary}[theorem]{Corollary}
\newtheorem{definition}[theorem]{Definition}
\newtheorem{remark}[theorem]{Remark}

% --- 記号の簡略化 ---
\newcommand{\Zp}{\mathbb{Z}_p}
\newcommand{\Zthree}{\mathbb{Z}^3}
\newcommand{\Sn}{S_n}
\newcommand{\Lfull}{Z_{\mathrm{full}}}
\newcommand{\Lzero}{Z_0}
\newcommand{\Rfilter}{R}
\newcommand{\rhoeff}{\rho_{\mathrm{eff}}}
\newcommand{\lpeak}{\lambda_{\mathrm{peak}}}

\doublespacing

\begin{document}

% ===== タイトル =====
\title{A Number-Theoretic Refinement of Stability Conditions\\
for Queues with Periodic Arrivals:\\
A Shell Filter Approach}

\author{Kenshirou Moriwaki\\
\small Independent Researcher, Nakashibetsu, Hokkaido, Japan\\
\small \texttt{contact@damrosch.net}}

\date{May 2026}
\maketitle

% ===== 要旨 =====
\begin{abstract}
This paper identifies a fundamental limitation of the classical queueing 
stability condition $\rho = \lambda/\mu < 1$ when applied to service systems 
with periodic arrival structures, and proposes a refined stability condition 
based on the shell filter limit theorem. For an arrival process with period 
governed by a prime $p$, we derive the necessary condition 
$\lambda_{\mathrm{peak}} < \mu(p-1)/p$, where $\lambda_{\mathrm{peak}}$ 
denotes the peak arrival rate within each period. This condition subsumes 
the classical condition as the special case $p \to \infty$ (Poisson arrivals). 
Applications to emergency medicine ($p=3$), tourism facilities ($p=7$), and 
urban transit ($p=5$) demonstrate that the refined condition correctly 
identifies instability in regimes where the classical condition erroneously 
predicts stability. Numerical verification confirms convergence of the shell 
filter ratio $R(s)$ to the theoretical value $(p-1)/p$.
\end{abstract}

\textbf{Keywords:} queueing theory; stability conditions; periodic arrivals; 
number theory; emergency medicine; shell filter

% ===== 第1節 =====
\section{Introduction}

\subsection{Background}

Japan's emergency medical system faces a severe congestion problem.
According to the Fire and Disaster Management Agency
\citep{FDMA2025}, emergency dispatches reached a record 7.72 million 
in 2024---the third consecutive year of record highs---with a national 
average hospital admission time of 44.6 minutes, approximately five 
minutes longer than the pre-pandemic level.

This crisis cannot be attributed solely to increasing patient volume. 
A more fundamental driver is the \emph{temporal heterogeneity of 
arrivals}: emergency transport demand exhibits clear periodicity across 
hours, days, and seasons. Morning outpatient peaks, nocturnal 
acute-episode clusters, and weekend exacerbations of chronic conditions 
combine to create periods of demand that far exceed capacity.

Yet the stability conditions widely used in hospital operations planning 
fail to capture this periodic structure.

\subsection{Limitations of Classical Theory}

The classical stability condition in queueing theory is
\begin{equation}\label{eq:classical}
  \rho = \frac{\lambda}{\mu} < 1,
\end{equation}
where $\lambda$ denotes the mean arrival rate and $\mu$ the mean service 
capacity \citep{Erlang1917,Kendall1953}. This condition underpins the 
$M/M/R$ model and has been applied extensively from telephone exchange 
design to hospital staffing.

\citet{Green2006} applied queueing theory to emergency department (ED) 
staffing and demonstrated that accounting for hourly variation in arrival 
rates can meaningfully reduce the rate of patients who leave without being 
seen (LWBS). However, even in that work, the stability condition itself 
remained the classical mean-based criterion.

\citet{Wiler2013} provided striking empirical evidence of the gap: 
a mere 10\% increase in the peak arrival rate above the observed mean 
caused the LWBS rate to jump from 3.9\% to 10.8\%---nearly a threefold 
increase. This occurs within a regime where condition~\eqref{eq:classical} 
still declares the system stable, revealing a structural blind spot.

The root cause is that the \emph{prime-periodic structure} of the arrival 
process is absent from the stability criterion. This paper provides a 
theoretically grounded answer from the framework of analytic number theory.

\subsection{Contributions}

The main contributions are as follows.
\begin{enumerate}
  \item \textbf{Theoretical refinement.} For an arrival process with prime 
    period $p$, we derive a necessary stability condition strictly tighter 
    than $\rho < 1$ (Section~3, Theorem~\ref{thm:refined}).
  \item \textbf{Unification.} The classical condition is recovered as the 
    limiting case $p \to \infty$ (Corollary~\ref{cor:unification}).
  \item \textbf{Applied validation.} In three-shift emergency medicine 
    ($p=3$) and weekly-cycle tourism ($p=7$), the refined condition 
    correctly identifies instability where the classical condition 
    predicts stability (Section~3.5).
\end{enumerate}

\subsection{Organization}

Section~2 introduces the shell filter in minimal form. Section~3 
develops the refined stability condition. Section~4 extends the framework 
and provides numerical verification. Section~5 concludes. Full proofs 
are in Appendix~B.

% ===== 第2節 =====
\section{Preliminaries: The Shell Filter}

\subsection{Periodic Classification of Arrivals}

Real service systems exhibit periodic demand structures. To describe this 
structure, we introduce the \emph{period parameter} $p$ (a prime number).

\begin{itemize}
  \item Three-shift scheduling (8 hours $\times$ 3): $p = 3$
  \item Weekly cycle (7 days): $p = 7$
  \item Five-day commute cycle: $p = 5$
\end{itemize}

\subsection{Definition of the Shell Filter}

For an observation window of length $n$, we represent an arrival pattern 
as a three-component vector $(x, y, z)$, where $x, y, z$ are integer-valued 
arrival intensities along independent channels (e.g., emergency, outpatient, 
referral).

\begin{definition}[L1 shell]
  The set of arrival patterns at observation length $n$ is
  \[
    S_n := \{(x,y,z) \in \mathbb{Z}^3 : |x|+|y|+|z| = n\},
    \quad |S_n| = 4n^2 + 2.
  \]
\end{definition}

\begin{definition}[Mod-$p$ filter count]
  Let $Q(x,y,z) = x^2+y^2+z^2$. Define
  \[
    N_0(n) := \#\{(x,y,z) \in S_n : Q(x,y,z) \equiv 0 \pmod{p}\},
  \]
  interpreted as the number of zero-load patterns at observation length $n$.
\end{definition}

\begin{definition}[Shell filter ratio]
  For $\mathrm{Re}(s) > 3/2$, define
  \begin{align*}
    Z_{\mathrm{full}}(s) &:= \sum_{n=1}^{\infty} |S_n| \cdot n^{-2s}
                          = 4\zeta(2s-2) + 2\zeta(2s),\\
    Z_0(s) &:= \sum_{n=1}^{\infty} N_0(n) \cdot n^{-2s},\\
    R(s)   &:= 1 - \frac{Z_0(s)}{Z_{\mathrm{full}}(s)}.
  \end{align*}
  $R(s)$ is the weighted proportion of arrival patterns that impose 
  effective load under the period-$p$ structure.
\end{definition}

% ===== 第3節 =====
\section{A Number-Theoretic Refinement of Stability Conditions}

\subsection{The Classical Condition and Its Limitation}

The classical stability condition $\rho = \lambda/\mu < 1$ embeds the 
assumption that arrivals follow a Poisson process. In emergency medicine, 
this assumption fails systematically: dispatch volumes differ severalfold 
between morning and midnight. When arrivals concentrate periodically, 
condition~\eqref{eq:classical} cannot detect chronic overload during 
specific windows even when the global average appears comfortable.

\subsection{Main Theorem and Its Queueing Interpretation}

\begin{theorem}[L1 shell mod-$p$ filter limit theorem; \citealt{Moriwaki2026b}]
\label{thm:main}
  For any odd prime $p$,
  \[
    \lim_{s \to (3/2)^+} R(s) = \frac{p-1}{p}.
  \]
\end{theorem}

The proof is given in Appendix~B. The queueing interpretation is as follows.

\begin{remark}[Load concentration guarantee]
  When the arrival process possesses a prime period $p$, the proportion 
  of time windows bearing effective load converges to $(p-1)/p$, and the 
  proportion of zero-load windows converges to $1/p$---a ratio guaranteed 
  by number theory, not approximation.
\end{remark}

\subsection{The Refined Stability Condition}

\begin{definition}[Effective utilization]
  Let $\lpeak$ denote the arrival rate during load-concentrated windows. 
  The \emph{effective utilization ratio} is
  \[
    \rhoeff := \frac{\lpeak \cdot p}{\mu \cdot (p-1)}.
  \]
\end{definition}

\begin{theorem}[Refined stability condition]
\label{thm:refined}
  A necessary condition for stability of a queueing system with prime 
  period $p$, strictly tighter than $\rho < 1$, is
  \[
    \rhoeff = \frac{\lpeak \cdot p}{\mu(p-1)} < 1,
    \quad \text{equivalently,} \quad
    \lpeak < \mu \cdot \frac{p-1}{p}.
  \]
\end{theorem}

\begin{proof}[Proof sketch]
  By Theorem~\ref{thm:main}, a fraction $(p-1)/p$ of the period bears 
  concentrated load. The effective demand in those windows is $\lpeak$, 
  against capacity $\mu$. Since $\lpeak \geq \lambda$ and $(p-1)/p < 1$, 
  the condition $\rhoeff < 1$ is strictly tighter than $\rho < 1$.
\end{proof}

\begin{corollary}[Recovery of the classical condition]
\label{cor:unification}
  The classical condition $\rho < 1$ is the special case $p \to \infty$.
  Since $(p-1)/p \to 1$ as $p \to \infty$, the refined condition reduces 
  to $\lpeak < \mu$, i.e., $\rho_{\mathrm{peak}} < 1$.
\end{corollary}

\subsection{Application to Emergency Medicine}

\citet{Wiler2013} show that a 10\% increase in peak arrival rate above 
the observed mean causes the LWBS rate to jump from 3.9\% to 10.8\%.
This occurs within a regime where $\rho < 1$, and represents precisely 
the structural instability that the refined condition detects.

\paragraph{Setting (based on \citealt{Green2006}).}
A three-shift ED ($p=3$, 8-hour shifts):
\begin{itemize}
  \item Mean arrival rate: $\lambda = 5.0$ patients/hour
  \item Service capacity: $\mu = 6.0$ patients/hour
  \item Peak arrival rate: $\lpeak = 10.0$ patients/hour
\end{itemize}

\paragraph{Classical condition:}
$\rho = 5.0/6.0 \approx 0.83 < 1$ \quad $\Rightarrow$ Stable (incorrect).

\paragraph{Refined condition ($p=3$):}
\[
  \rhoeff = \frac{10.0 \times 3}{6.0 \times 2} = \frac{30.0}{12.0} 
          = 2.50 > 1 \quad \Rightarrow \text{Unstable (correct).}
\]

The gap between these two assessments provides a number-theoretic 
explanation for the documented LWBS spike \citep{Wiler2013}.

% ===== 第4節 =====
\section{Extensions and Numerical Verification}

\subsection{Numerical Convergence of $R(s)$}

Table~\ref{tab:convergence} reports the finite-truncation approximation 
$R_N(s)$ for $N=80$ as $s \to (3/2)^+$.

\begin{table}[h]
\centering
\caption{Convergence of $R_N(s)$ as $s \to (3/2)^+$ ($N = 80$)}
\label{tab:convergence}
\begin{tabular}{ccccc}
\toprule
$s$ & $p=3$ & $p=5$ & $p=7$ \\
\midrule
1.550 & 0.8095 & 0.8538 & 0.9349 \\
1.520 & 0.7981 & 0.8478 & 0.9296 \\
1.510 & 0.7943 & 0.8458 & 0.9278 \\
1.505 & 0.7924 & 0.8448 & 0.9269 \\
\midrule
Theory & \textbf{0.6667} & \textbf{0.8000} & \textbf{0.8571} \\
\bottomrule
\end{tabular}
\end{table}

The overshoot is a known artifact of the non-commutativity of the double 
limit ($N \to \infty$ and $s \to (3/2)^+$), documented in \citet{Moriwaki2026a}.

\subsection{Refined Stability Thresholds by Prime Period}

\begin{table}[h]
\centering
\caption{Refined stability thresholds by prime period ($\mu = 6.0$)}
\label{tab:thresholds}
\begin{tabular}{cccc}
\toprule
Period $p$ & $(p-1)/p$ & Peak threshold $\lpeak$ & Ratio to classical \\
\midrule
3  & 0.6667 & 4.00 & 0.667 \\
5  & 0.8000 & 4.80 & 0.800 \\
7  & 0.8571 & 5.14 & 0.857 \\
11 & 0.9091 & 5.45 & 0.909 \\
13 & 0.9231 & 5.54 & 0.923 \\
$\infty$ & 1.0000 & 6.00 & 1.000 \\
\bottomrule
\end{tabular}
\end{table}

\subsection{Application to Tourism and Urban Transit}

\begin{table}[h]
\centering
\caption{Classical vs.\ refined conditions across three domains}
\label{tab:comparison}
\begin{tabular}{lccccc}
\toprule
Domain & $p$ & Classical $\rho$ & & Refined $\rhoeff$ & \\
\midrule
Emergency medicine (3-shift) & 3 & 0.83 & Stable$^*$ & 2.50 & Unstable \\
Tourism (weekly)              & 7 & 0.76 & Stable$^*$ & 1.87 & Unstable \\
Urban transit (commute)       & 5 & 0.75 & Stable$^*$ & 1.75 & Unstable \\
\bottomrule
\multicolumn{6}{l}{$^*$ Incorrect assessment.}
\end{tabular}
\end{table}

Across all three domains, the classical condition produces false-stable 
assessments while the refined condition correctly identifies instability.

% ===== 第5節 =====
\section{Conclusion}

This paper derived a refined stability condition for queueing systems 
with periodic arrival structures based on the shell filter limit theorem.
The key result is that when the arrival period is governed by a prime $p$,
the permissible peak arrival rate satisfies $\lpeak < \mu(p-1)/p$---strictly
tighter than the classical condition, which is recovered as $p \to \infty$.

Applications to emergency medicine, tourism, and urban transit demonstrate 
that the refined condition correctly identifies structural instability in 
regimes where the classical condition fails.

\paragraph{Future directions.}
\begin{enumerate}
  \item Estimation of $p$ from observed arrival data (e.g., FDMA dispatch records).
  \item Extension to composite periods via the Chinese Remainder Theorem.
  \item Higher-dimensional generalization: $d$-channel arrivals and 
    $\lim R(s) = (p^{d-2}-1)/p^{d-2}$.
\end{enumerate}

% ===== 参考文献 =====
\bibliographystyle{plainnat}
\begin{thebibliography}{9}

\bibitem[Erlang(1917)]{Erlang1917}
Erlang, A.~K. (1917).
Solution of some problems in the theory of probabilities of significance 
in automatic telephone exchanges.
\emph{Post Office Electrical Engineers' Journal}, 10, 189--197.

\bibitem[Fire and Disaster Management Agency(2025)]{FDMA2025}
Fire and Disaster Management Agency, Japan (2025).
\emph{Status of Emergency and Rescue Operations, 2025 Edition}.
Ministry of Internal Affairs and Communications.

\bibitem[Green et al.(2006)]{Green2006}
Green, L.~V., Soares, J., Giglio, J.~F., \& Green, R.~A. (2006).
Using queueing theory to increase the effectiveness of emergency 
department provider staffing.
\emph{Academic Emergency Medicine}, 13(1), 61--68.

\bibitem[Kendall(1953)]{Kendall1953}
Kendall, D.~G. (1953).
Stochastic processes occurring in the theory of queues and their 
analysis by the method of the imbedded Markov chain.
\emph{The Annals of Mathematical Statistics}, 24(3), 338--354.

\bibitem[Moriwaki(2026a)]{Moriwaki2026a}
Moriwaki, K. (2026a).
Density of Eisenstein primes on hexagonal shells: A numerical 
verification of Chebotarev's theorem via Ces\`{a}ro averaging.
\emph{damrosch.net}.

\bibitem[Moriwaki and Claude(2026b)]{Moriwaki2026b}
Moriwaki, K., \& Claude (Anthropic). (2026b).
L1 shell mod $p$ filter limit theorem.
\emph{Unpublished manuscript}.

\bibitem[Wiler et al.(2013)]{Wiler2013}
Wiler, J.~L., Bolandifar, E., Griffey, R.~T., Poirier, R.~F., \& Olsen, T. (2013).
An emergency department patient flow model based on queueing theory principles.
\emph{Academic Emergency Medicine}, 20(9), 939--946.

\end{thebibliography}

% ===== 付録B =====
\appendix
\section*{Appendix B: Proof of Theorem~\ref{thm:main}}
\addcontentsline{toc}{section}{Appendix B: Proof of Theorem 1}

\subsection*{B.1 Notation and the $\zeta$-Representation of $Z_{\mathrm{full}}$}

\begin{lemma}[$\zeta$-representation of $Z_{\mathrm{full}}$]
\label{lem:zfull}
  For $\mathrm{Re}(s) > 3/2$,
  \[
    Z_{\mathrm{full}}(s) 
    = \sum_{n=1}^{\infty}(4n^2+2)\,n^{-2s}
    = 4\zeta(2s-2) + 2\zeta(2s).
  \]
\end{lemma}

\begin{proof}
  Substituting $|S_n|=4n^2+2$:
  \[
    Z_{\mathrm{full}}(s) 
    = 4\sum_{n=1}^{\infty}n^{2-2s} + 2\sum_{n=1}^{\infty}n^{-2s}
    = 4\zeta(2s-2)+2\zeta(2s).
  \]
  The first series converges for $\mathrm{Re}(2s-2)>1$, i.e., 
  $\mathrm{Re}(s)>3/2$; the second for $\mathrm{Re}(s)>1/2$.
  Hence the convergence domain is $\mathrm{Re}(s)>3/2$.
\end{proof}

\subsection*{B.2 Key Lemma}

\begin{lemma}
\label{lem:key}
  For any odd prime $p$,
  \[
    \#\{(x,y,z)\in\mathbb{Z}_p^3 : x^2+y^2+z^2\equiv 0\pmod{p}\} = p^2.
  \]
\end{lemma}

\begin{proof}
  Solving for $z$: the condition is $z^2\equiv-(x^2+y^2)\pmod{p}$.
  By the Legendre symbol formula $\#\{z\in\mathbb{Z}_p:z^2\equiv c\}=1+(c/p)_L$,
  \[
    N_0 = \sum_{x,y\in\mathbb{Z}_p}\!\!\left[1+\left(\frac{-(x^2+y^2)}{p}\right)_{\!L}\right]
        = p^2 + \left(\frac{-1}{p}\right)_{\!L} T,
  \]
  where $T=\sum_{x,y\in\mathbb{Z}_p}((x^2+y^2)/p)_L$.
  We show $T=0$. For $x=0$: $\sum_y(y^2/p)_L=p-1$.
  For $x\neq 0$ (fixing $c=x^2\neq 0$), the substitution
  $v=y^2/c$, $w=1-1/v$ yields $\sum_{v\neq 0,1}(v(v-1)/p)_L=-1$.
  Hence $T=(p-1)+(p-1)(-1)=0$, giving $N_0=p^2$.
\end{proof}

\subsection*{B.3 Coefficient Sum Lemma}

\begin{lemma}
\label{lem:coeff}
  Writing $N_0(n)=A_r q^2+O(q)$ for $n=pq+r$,
  we have $\sum_{r=0}^{p-1}A_r = 4p^2$.
\end{lemma}

\begin{proof}[Proof sketch]
  Expanding via exponential sums and using Lemma~\ref{lem:key}, 
  the contributions from $t\neq 0$ cancel, giving 
  $\sum_r A_r=(1/p)\cdot p\cdot 4p^2=4p^2$.
\end{proof}

\subsection*{B.4 Proof of Theorem~\ref{thm:main}}

\begin{proof}
  The leading term of $Z_0(s)$ as $s\to(3/2)^+$ satisfies
  \[
    Z_0(s) \sim \frac{\sum_r A_r}{p^{2s}}\,\zeta(2s-2)
            = \frac{4p^2}{p^{2s}}\,\zeta(2s-2)
            \sim \frac{4}{p}\,\zeta(2s-2),
  \]
  since $p^{2s}\to p^3$. By Lemma~\ref{lem:zfull},
  $Z_{\mathrm{full}}(s)\sim 4\zeta(2s-2)$. Therefore
  \[
    \frac{Z_0(s)}{Z_{\mathrm{full}}(s)} \to \frac{4/p}{4} = \frac{1}{p},
    \quad\text{and}\quad
    R(s) = 1 - \frac{Z_0}{Z_{\mathrm{full}}} \to \frac{p-1}{p}.\qed
  \]
\end{proof}

\end{document}
